\documentclass[../readings.tex]{subfiles}

\begin{document}
\subsection{Koopman Theory, DMD, and SINDy}
\label{sub:koopman}

Data-driven discovery of linear models for nonlinear dynamics.

\subsubsection{Koopman Operator Theory}

\addDef{Nonlinear Dynamics and Observables}{%
$\bx^{(k+1)} = \mathbf{F}(\bx^{(k)})$ on $\mathcal{M} \subseteq \fR^n$.
\begin{itemize}
    \item \textbf{Observable:}\enspace Scalar function $g: \mathcal{M} \to \fR$.
    \item \textbf{Koopman Operator $\mathcal{K}$:}\enspace Linear operator acting on observables: $(\mathcal{K}g)(\bx) = g(\mathbf{F}(\bx))$.
\end{itemize}
\textbf{Lifting:}\enspace $\mathcal{K}$ is linear even if $\mathbf{F}$ is nonlinear, at the cost of being infinite-dimensional.
}

\addDef{Koopman Eigenfunctions}{%
Observable $\varphi$ satisfying $\mathcal{K}\varphi = \mu \varphi$.
\begin{itemize}
    \item \textbf{Temporal Evolution:}\enspace $\varphi(\bx^{(k)}) = \mu^k \varphi(\bx^{(0)})$.
    \item \textbf{Significance:}\enspace Provides a global linear coordinate system for nonlinear dynamics.
\end{itemize}
}

\subsubsection{Dynamic Mode Decomposition (DMD)}

\addDef{DMD Algorithm}{%
Finds the best linear fit $\bY \approx \bA\bX$ for snapshot matrices:
\begin{align*}
    \bX = [\bx^{(0)} \dots \bx^{(m-1)}], \quad \bY = [\bx^{(1)} \dots \bx^{(m)}].
\end{align*}
\begin{enumerate}
    \item Compute truncated SVD: $\bX \approx \bU_r \bsigma_r \bV_r^T$.
    \item Construct $r \times r$ reduced operator: $\tilde{\bA} = \bU_r^T \bY \bV_r \bsigma_r^{-1}$.
    \item DMD eigenvalues $\{\lambda_i\}$ and modes $\{\boldsymbol{\phi}_i\}$ are recovered from $\tilde{\bA}$.
\end{enumerate}
}

\addNote{%
\textbf{Interpretation:}\enspace DMD decomposes data into spatiotemporal modes $\boldsymbol{\phi}_i$ with associated frequencies and growth rates. It is the data-driven approximation of the Koopman operator using the identity dictionary.
}

\subsubsection{SINDy: Sparse Identification}

\addDef{Sparse Regression for Dynamics}{%
Identifies governing equations $\dot{\bx} = \mathbf{\Xi}^T \boldsymbol{\Theta}(\bx)$ where $\boldsymbol{\Theta}$ is a library of candidate functions (monomials, sin, etc.).
\begin{itemize}
    \item \textbf{STLS Algorithm:}\enspace Sequential Thresholded Least Squares. Iteratively solves least squares and thresholds small coefficients to promote sparsity.
    \item \textbf{Outcome:}\enspace Parsimonious, interpretable models (e.g., recovering Lorenz equations).
\end{itemize}
}

\subsubsection{Exercises}

\addExcr{%
\begin{enumerate}
    \item Prove the linearity of the Koopman operator.
    \item Implement DMD for a 2D damped oscillator. Verify the identified eigenvalues.
    \item Build a polynomial library $\boldsymbol{\Theta}$ for $\bx \in \fR^2$ up to degree 2.
    \item Use SINDy (via STLS) to identify the Lotka-Volterra predator-prey model from noisy data.
\end{enumerate}
}

\addPrf{%
For any state $\bx \in \mathcal{M}$, observables $g, h$, and scalars $\alpha, \beta$, consider the action of $\mathcal{K}$ on the linear combination:
\begin{align*}
    [\mathcal{K}(\alpha g + \beta h)](\bx) &= (\alpha g + \beta h)(\mathbf{F}(\bx)) \\
    &= \alpha g(\mathbf{F}(\bx)) + \beta h(\mathbf{F}(\bx)) \\
    &= \alpha (\mathcal{K}g)(\bx) + \beta (\mathcal{K}h)(\bx) \\
    &= [\alpha \mathcal{K}g + \beta \mathcal{K}h](\bx).
\end{align*}
Since this holds for all $\bx$, $\mathcal{K}(\alpha g + \beta h) = \alpha \mathcal{K}g + \beta \mathcal{K}h$.
}

\end{document}
